\title{DeepSeekMath: Pushing the Limits of Mathematical Reasoning in Open Language Models}
The resulting model, DeepSeekMath-Instruct 7B, surpasses all existing 7B-scale models and demonstrates competitive performance compared to 70B open-source models that have been instruction-tuned. PPO seeks to optimize LLMs by maximizing the following clipped surrogate objective:

\begin{equation}
\footnotesize
                \mathcal{J}_{PPO}(\theta) = \mathbb{E}{[q \sim P(Q), o \sim \pi_{\theta_{old}}(O|q)]} \frac{1}{|o|} \sum_{t=1}^{|o|} \min \left[ \frac{\pi_\theta(o_{t} | q, o_{<t})}{\pi_{\theta_{old}}(o_{t} | q, o_{<t})} A_{t}, \text{clip} \left( \frac{\pi_\theta(o_{t} | q, o_{<t})}{\pi_{\theta_{old}}(o_{t} | q, o_{<t})}, 1 - \varepsilon, 1 + \varepsilon \right)  A_{t} \right] ,
\end{equation}

Here, $\pi_{\theta}$ denotes the current policy, while $\pi_{\theta_{old}}$ refers to the prior policy model. For each mathematical dataset, we individually fine-tune the model using 150B training tokens. This dataset comprises 776K training samples in total. This trend is reasonable since, during the initial phase, the actor and the SFT model remain closely aligned, resulting in only minimal variations in sampled outputs. For each domain, we determine the fraction of web pages retrieved in the initial pass. In outcome supervision, this normalized reward is delivered as a final scalar at the end of each output $o_i$, and each token within $o_i$ receives the same advantage value, specifically, $\hat{A}_{i, t} = \widetilde{r}_i = \frac{r_i- {\rm mean}(\mathbf{r})}{{\rm std}(\mathbf{r})}$. Our experiments divide the reward mechanism into `Rule' and `Model' categories. \item \textbf{PPO/GRPO}: These methods initialize the policy model with SFT weights and subsequently optimize it through reinforcement learning using outputs sampled dynamically from the evolving policy. The policy model undergoes only one parameter update per exploration round. We leverage the open-source implementation of fastText\footnote{\url{https://fasttext.cc}} with the following configuration: a vector dimensionality of 256, learning rate set to 0.1, word n-grams up to length 3, a minimum word frequency threshold of 3, and training run for 3 epochs. \end{itemize}

These models serve a general-purpose function and have typically been subjected to multiple stages of alignment procedures. Because the value estimator in PPO generally has a similar parameter scale as the policy, this leads to significant computational and memory costs. After repeating this collection process for four rounds, our dataset reaches 35.5M math-related web pages, amounting to 120B tokens. The model’s parameters are initialized with DeepSeek-Coder-Base-v1.5 7B \citep{deepseek-coder}, followed by continued pretraining over 500B tokens. However, the blended one-stage training with both code and math tokens appears to hamper performance on non-tool math reasoning tasks. Despite strong quantitative reasoning performance, DeepSeekMath~lags behind closed models when it comes to geometry and theorem-proving tasks. Furthermore, in contrast to PPO, which incorporates the KL penalty within the reward function, GRPO regularizes by directly including a KL divergence loss between the updated and reference policies, thereby simplifying the calculation for $\hat{A}_{i,t}$. \end{itemize}

For our analyses, DeepSeek-LLM 1.3B was trained for 150B tokens, and DeepSeek-Coder-Base-v1.5 7B for 40B tokens, each time using an isolated arXiv-based dataset. \end{abstract}

\section{Introduction}

The advent of large language models (LLMs) has dramatically transformed AI-driven mathematical reasoning, leading to marked progress on quantitative benchmarks such as \citep{MATH} and geometry challenges exemplified by \citep{trinh2024solving}. The answer is determined by executing the generated program. For experiments with Proof-Pile-2, we adopt the arXiv:Web:Code proportion of 2:4:1 as prescribed in \cite{llemma}. Nevertheless, there is no guarantee that the reward signal is consistently dependable, particularly in highly intricate problem domains. \paragraph{Stepwise Mathematical Problem Solving}

The model’s mathematical problem-solving abilities are assessed using few-shot chain-of-thought prompting \citep{cot} on eight separate English and Chinese benchmarks. These step-wise rewards are normalized using the group mean and standard deviation: $\widetilde{r}_i^{index(j)} = \frac{r_i^{index(j)} - {\rm mean(\mathbf{R})}}{{\rm std(\mathbf{R})}}$. The learning rate is managed via a multi-step scheduling approach: starting with a 2,000-step warmup to its maximum, followed by a reduction to 31.6\% after 80\% of training, and a final decrease to 10.0\% at 90\% completion. Within reinforcement learning, this typically takes the form of a neural reward model. \subsubsection{How to Achieve More Effective RL?}
\label{sec:effec_rl}
Our results confirm that RL can be highly effective in mathematical reasoning scenarios. \item Our results indicate that GRPO yields substantial improvements in the instruction-tuned model DeepSeekMath-Instruct when leveraging only instruction-tuning data. Equations \ref{eq:GC-RFT} and \ref{eq:GC-GRPO} delineate an essential distinction: GRPO uniquely calibrates its gradient coefficients by considering the reward model’s values, enabling tailored reinforcement or penalization of responses in proportion to their assessed merit. Following this, the reference model is synchronized with the updated policy model, and further policy optimization is performed using the refreshed reward model. Several factors remain unaddressed, including:

\begin{itemize}[topsep=0pt]
    \item The contribution of arXiv-derived tokens to specific mathematical tasks not considered here—such as converting formal mathematical statements or proofs to their informal counterparts (informalization of theorems);
    \item Potential interactions between arXiv text and other data modalities when included jointly in pre-training;
    \item Possible advantages from arXiv data that could surface only at much larger model scales. Pre-training with mathematical data enhances both comprehension and reasoning capacities of the model. All training is performed within the HAI-LLM \citep{haillm} system, chosen for its efficiency and low resource demands. GRPO eliminates the need for a critic network by utilizing group scores to compute the baseline, which drastically lowers computational requirements in comparison with Proximal Policy Optimization (PPO). Each problem is matched with a solution presented in one of several formats: chain-of-thought (CoT) \citep{cot}, program-of-thought (PoT) \citep{pot,pal}, or reasoning incorporating external tools \citep{tora}. A likely explanation is that the capacity limitations of DeepSeek-LLM 1.3B may constrain the simultaneous learning of both coding and mathematical patterns. In its initial stage, the classifier is trained with examples from OpenWebMath \citep{openwebmath} as positive samples, while various unrelated web pages constitute the negative set. \paragraph{Reflection on Data Source} We categorize data sources into two principal types: online sampling and offline sampling. We also intend to pursue the research directions highlighted in Section \ref{sec:effec_rl} to achieve more effective reinforcement learning strategies for large language models. However, in the case of LLMs, the reward model typically assesses only the final token of a sequence, presenting challenges for training a value function capable of yielding precise, per-token estimates. \item \textbf{Direct Preference Optimization (DPO)}: DPO advances SFT-trained models by further fine-tuning on paired candidate responses drawn from the model itself, optimizing with a preference-based DPO loss function. \item \textbf{Rejection Sampling Fine-tuning (RFT)}: RFT continues model adaptation using a subset of model-generated responses, filtered for correctness based on human-assigned SFT prompts. A fastText model \citep{joulin2016fasttext} is subsequently trained to retrieve further web pages similar in content to OpenWebMath. \end{itemize}

\item \textbf{Formal Mathematics}: DeepSeekMath-Base is assessed on the informal-to-formal theorem proving benchmark from \citep{dsp_proof}, utilizing the miniF2F dataset \citep{minif2f} and employing Isabelle \citep{isabelle} as the proof assistant. To gain a deeper understanding of the mechanisms behind RL's effectiveness, we assess both Pass@K and Maj@K accuracies of the Instruct and RL models on two benchmarks. GRPO's group-relative advantage formulation complements the pairwise or comparative structure commonly used in reward modeling, as reward models are typically developed on comparison datasets involving multiple responses per question. In the two-stage configuration, the initial phase of code training produces measurable progress, and this momentum persists and amplifies during the math training phase, resulting in the overall highest scores. As detailed in Table \ref{tab:base_math_cot}, DeepSeekMath-Base 7B achieves top performance among open-source base models—including Mistral 7B \citep{mistral}, a widely adopted general-purpose model, and the recently introduced Llemma 34B \citep{llemma}, which has been further trained on mathematical data from Proof-Pile-2 \citep{llemma}—across all eight benchmark datasets. Models of smaller size, when trained on data of superior quality, can also deliver competitive outcomes. GRPO removes the need for a critic network and instead computes a baseline using group-based reward statistics, thereby greatly reducing computational requirements. We test DeepSeekMath-Base 7B on the informal-to-formal proof translation task from \citep{dsp_proof}, where the objective is to convert an informal mathematical claim, along with its formal version and informal proof, into a corresponding formal proof. Conversely, most other math corpora focus on English content and consequently fail to enhance, or may even degrade, performance for Chinese mathematical problems. \textbf{Math Pre-Training at Scale}

\begin{itemize}[topsep=0pt]
    \item This study demonstrates that Common Crawl, an openly available data source, possesses substantial value for mathematical applications. By the fourth iteration, we observe that approximately 98\% of the web pages had already been acquired by the third cycle, prompting us to halt further collection. The `Rule' type assesses response quality solely on answer correctness, whereas the `Model' type involves training a reward model—using data labeled via the rule-based criteria—to evaluate responses. Using the trained fastText model, we extract mathematical pages from this deduplicated set. Consistently, we observed little to no benefit—and in some cases performance degradation—across a range of mathematical reasoning assessments, regardless of the models’ scale. \paragraph{Data Source} The data source underpins all training methodologies. \item \textbf{Online Rejection Sampling Fine-tuning (Online RFT)}: Distinct from the traditional RFT, Online RFT reinitializes policy learning from the SFT checkpoint, updating parameters through real-time sampling of policy outputs. In the GRPO framework, the policy model employs a learning rate of 1e-6, while the KL coefficient is maintained at 0.04. DeepSeekMath~7B attains a notable 51.7\% score on the MATH benchmark at competition level, doing so independently of external tool support or ensemble voting methods, thereby nearing the performance of Gemini-Ultra and GPT-4. Evaluation of DeepSeekMath-RL 7B is performed using the same set of benchmarks as DeepSeekMath-Instruct 7B. \end{itemize}

\subsection{Training and Evaluating DeepSeekMath-Instruct 7B}

Here, we detail the development of DeepSeekMath-Instruct 7B, which is instruction-tuned for mathematical reasoning using the DeepSeekMath-Base model as its starting point. Notably, in our preliminary assessments, the model fails to solve questions related to triangles and ellipses, which may reflect biases introduced during data selection in pre-training and fine-tuning. \item \textbf{Open-source models} consist of: (1) DeepSeek-LLM-Chat 67B \citep{deepseek-llm}, (2) Qwen 72B \citep{qwen}, (3) SeaLLM-v2 7B \citep{seallm}, and (4) ChatGLM3 6B \citep{chatglm3}, all of which are generalist models, along with those designed with additional mathematical capabilities: (5) InternLM2-Math 20B~\footnote{\url{https://github.com/InternLM/InternLM-Math}}, an extension of InternLM2 that undergoes math-specific training and subsequent instruction tuning; (6) Math-Shepherd-Mistral 7B, which applies PPO-based training \citep{schulman2017proximal} to Mistral 7B \citep{mistral} using a process-supervised reward mechanism; (7) the WizardMath models \citep{wizardmath}, which refine mathematical reasoning for Mistral 7B and Llama-2 70B \citep{llama2} via evolve-instruct (an instruction tuning strategy using machine-generated instructions) alongside PPO and training tasks primarily sourced from GSM8K and MATH; (8) MetaMath 70B \citep{metamath}, which fine-tunes Llama-2 70B on an enriched GSM8K and MATH corpus; (9) ToRA 34B \cite{tora}, a version of CodeLlama 34B adapted for tool-augmented mathematical reasoning through fine-tuning; (10) MAmmoTH 70B \citep{MathInstruct}, which is Llama-2 70B fine-tuned on MathInstruct following an instruction-tuning procedure. Motivated by these limitations, our focus is on investigating reinforcement learning algorithms that maintain robustness in the face of noisy or unreliable reward signals. On additional evaluation sets, it remains highly competitive with DeepSeek-LLM-Chat 67B, which previously led the field despite being ten times larger. \section{Supervised Fine-Tuning}

\subsection{SFT Data Curation}

We assemble an instruction-tuning dataset for mathematics that encompasses English and Chinese problems from a range of mathematical disciplines and varying degrees of difficulty. Nonetheless, it is important to acknowledge limitations in this interpretation. \item \textbf{General Training for 400B Tokens $\rightarrow$ Math Training for 150B Tokens}: As a baseline for comparison, we replace the initial code tokens with general tokens (drawn from DeepSeek-AI’s comprehensive general corpus) in the first pretraining phase. To decrease the scale of Common Crawl, we apply both URL deduplication and near-deduplication strategies, ultimately filtering the dataset to 40 billion HTML pages. Additionally, we propose Group Relative Policy Optimization (GRPO)—an adaptation of Proximal Policy Optimization (PPO)—which enhances mathematical reasoning skills while lowering memory requirements. \item \textbf{Chinese mathematical datasets}: We gather K-12 level mathematical questions in Chinese, covering 76 different sub-domains such as linear equations, with each item accompanied by CoT and tool-based reasoning solutions. \paragraph{Reflection on Gradient Coefficient} The algorithms use the processed reward signal to determine the gradient coefficient, which governs the update to the model parameters. While DeepSeekMath-Instruct's results are on par with proprietary Chinese models like GLM-4 and Baichuan-3 on MATH, it does not yet match the accuracy of GPT-4 or Gemini Ultra. The reward model is then continually refined, leveraging a replay strategy that retains 10\% of prior samples. Comparative evaluation is performed against prominent contemporaneous models:

\begin{itemize}[topsep=0pt]
    \item \textbf{Closed-source models} evaluated include: (1) the GPT model family—especially GPT-4 \citep{gpt4} and GPT-4 Code Interpreter~\footnote{\url{https://openai.com/blog/chatgpt-plugins\#\#code-interpreter}}, (2) Gemini Ultra and Pro \citep{gemini}, (3) Inflection-2 \citep{inflection-2}, (4) Grok-1~\footnote{\url{https://x.ai/model-card}}, in addition to Chinese-developed models such as (5) Baichuan-3~\footnote{\url{https://www.baichuan-ai.com}}, (6) the latest entry from the GLM family, GLM-4~\footnote{\url{https://open.bigmodel.cn/dev/api\#glm-4}} \citep{glm}. As demonstrated in Figure \ref{fig:maj-pass}, RL elevates Maj@K but does not impact Pass@K, suggesting that the strength of RL lies in generating a more reliable output distribution. Our comprehensive ablation studies demonstrate that web pages are a promising and rich resource for obtaining high-quality mathematical datasets, whereas arXiv content does not contribute as much as anticipated. Our approach relies on constructing the DeepSeekMath Corpus, a comprehensive and high-quality pre-training dataset containing 120B math tokens. Both RFT and DPO are classified as using the offline approach, while Online RFT and GRPO leverage the online approach. On English benchmarks, DeepSeekMath-Base matches the performance of the proprietary Minerva 540B \citep{minerva}, and notably outperforms all available open-source base models (such as Mistral 7B \citep{mistral} and Llemma-34B \citep{llemma}), whether or not those have been math-pretrained, often by a large margin. Following this, we employ the classifier to extract further positive samples from CC; these are subjected to manual annotation for greater precision. As an initial seed, OpenWebMath \citep{openwebmath}—a curated collection of high-quality mathematics-related texts sourced from the web—is selected. Table \ref{tab:sft_rl_math} presents comparative results for both open-source and proprietary models, incorporating chain-of-thought as well as tool-assisted reasoning approaches, evaluated on English and Chinese test sets. Our assessments show that this extensive corpus supports superior model performance: DeepSeekMath-Base 7B attains 64.2\% on GSM8K \citep{gsm8k} and 36.2\% on the competition-grade MATH dataset \citep{MATH}, exceeding the results of Minerva 540B \citep{minerva}. \end{itemize}

\paragraph{Results}
The evaluation results for each training approach are summarized in Table \ref{tab:code-math} and Table \ref{tab:code-math-general-eval}. To address this, we augment our seed dataset by identifying further mathematics-focused web domains, thereby enhancing the fastText model's effectiveness. We further propose the Group Relative Policy Optimization (GRPO) method, a modified reinforcement learning algorithm derived from Proximal Policy Optimization (PPO) \citep{schulman2017proximal}. This performance edge holds even over much larger models (e.g., Qwen 72B) or those that have benefited from specialized reinforcement learning tailored for mathematical tasks (such as WizardMath-v1.1 7B). \paragraph{Natural Language Understanding, Reasoning, and Code}

For natural language understanding, we assess the model on MMLU \citep{mmlu}, for reasoning on BBH \citep{bbh}, and for programming proficiency using HumanEval \citep{codex} and MBPP \citep{mbpp}. \subsubsection{From PPO to GRPO} Proximal Policy Optimization (PPO) \citep{schulman2017proximal} is a well-established actor-critic reinforcement learning technique that is commonly employed for further tuning of LLMs during RL stages \citep{ouyang2022training}. Specifically, for a prompt $q$ and corresponding samples $\{o_1, o_2, \cdots, o_G\}$, a process-based reward model delivers a sequence of scores for each step: $\mathbf{R} = \{ \{r_1^{index(1)},\cdots,r_1^{index(K_1)}\}, \cdots,  \{r_G^{index(1)},\cdots,r_G^{index(K_G)}\} \}$, where $index(j)$ locates the terminal token of the $j$-th step, and $K_i$ indicates the number of steps in output $o_i$. To systematically assess the influence of code training on mathematical reasoning, we investigated two methodological variations: a two-stage training protocol and a single-stage training approach. Looking ahead, we aim to refine our data selection pipeline further, curating higher-quality pre-training corpora. Reward model training data are prepared in line with the methodology in \citep{wang2023math}. The scalar $\varepsilon$ functions as a clipping hyper-parameter as introduced in PPO, contributing to the robustness of the optimization process. Moreover, when code and math tokens are merged in a one-stage training procedure, the model not only sidesteps the catastrophic forgetting observed with sequential fine-tuning but also achieves improvements in both coding skills (see Table \ref{tab:code-math-general-eval}) and math tasks assisted by program execution (see Table \ref{tab:code-math}). In both sequential and joint training paradigms, pretraining with code appears to significantly enhance program-assisted mathematical problem-solving. \textbf{Two-Stage Training}

\begin{itemize}[topsep=0pt]
    \item \textbf{Code Training for 400B Tokens $\rightarrow$ Math Training for 150B Tokens}: DeepSeek-LLM 1.3B is first pretrained with 400B tokens of code, followed by a further 150B tokens of mathematics-specific training. \paragraph{Reward Function} The reward function constitutes the primary origin of training signals. We systematically evaluated models of differing capacities—specifically, DeepSeek-LLM 1.3B and DeepSeek-Coder-Base-v1.5 7B \citep{deepseek-coder}—using arXiv text drawn from distinct preprocessing regimes:

\begin{itemize}[topsep=0pt]
    \item \textbf{MathPile} \citep{mathpile}: a corpus comprising 8.9B tokens, primarily consisting (over 85\%) of scientific arXiv material, and subject to cleaning and heuristic filtering;
    \item \textbf{ArXiv-RedPajama} \citep{redpajama}: all arXiv LaTeX source files with non-essential components—such as preambles, comments, macros, and bibliographies—excised, yielding a total of 28.0B tokens. \subsubsection{Evaluation Results}

\textbf{The DeepSeekMath~Corpus exhibits superior quality, includes a breadth of multilingual mathematical material, and is unmatched in scale.}

\begin{itemize}[topsep=0pt]
    \item \textbf{High-quality}:
    To assess the downstream capabilities, we test performance on 8 mathematical benchmarks, utilizing few-shot chain-of-thought prompting \cite{cot}. Importantly, this approach can be transferred to additional specialized domains, including but not limited to code. The English dataset provides broad mathematical domain coverage, including but not limited to algebra, probability, number theory, calculus, and geometry. \EndFor 
  \end{algorithmic}
  \textbf{Output} $\pi_\theta$
  \label{alg:iter-grpo}
\end{algorithm}

\subsubsection{Outcome Supervision RL with GRPO} More formally, for a given question $q$, the procedure draws a set of responses $\{o_1, o_2, \cdots, o_G\}$ from the prior version of the policy $\pi_{\theta_{old}}$. For each task, formal proofs in Isabelle are generated with few-shot prompts. Importantly, references to the DeepSeekMath Corpus here pertain to the second version of our dataset, comprising 89 billion tokens. \paragraph{Mathematical Problem Solving with Tool Use} For assessing program-supported mathematical reasoning, we use few-shot program-of-thought prompting \citep{pot,pal} on the GSM8K and MATH datasets. \end{itemize}

\subsubsection{Training Setting}
\label{sec:quality-policy}
We conduct math-specific fine-tuning on a general-purpose language model comprising 1.3B parameters, utilizing the same architecture as the DeepSeek LLMs \citep{deepseek-llm}, henceforth referred to as DeepSeek-LLM 1.3B. Illustrated in Figure \ref{fig:data_collect}, an iterative procedure is employed to expand a large-scale mathematical dataset, originating from a small, high-quality seed collection—such as a specialized math dataset. Our methodology involves partitioning the entire Common Crawl dataset into non-overlapping domains, where each domain comprises web pages with the identical base URL. In line with the methodology of \cite{dsp_proof}, our approach generates proof outlines that are subsequently elaborated by the automated prover Sledgehammer \citep{sledgehammer}. \subsubsection{ArXiv Papers Appear Unhelpful for Enhancing Mathematical Reasoning}

Incorporating arXiv papers into mathematical pre-training datasets has become standard practice \citep{minerva,gpt-f,llemma,mathpile}. This pattern persisted across benchmarks targeting quantitative reasoning (e.g., GSM8K and MATH; see Table \ref{tab:arxiv-cot}), multiple-choice tasks like MMLU-STEM (Table \ref{tab:arxiv-cot}), and formal mathematics benchmarks, such as miniF2F (Table \ref{tab:arxiv_atp}). Following the pre-training stage, we perform mathematical instruction tuning on DeepSeekMath-Base, utilizing datasets featuring chain-of-thought \citep{cot}, program-of-thought \citep{pot,pal}, and tool-integrated reasoning \citep{tora} paradigms. Additionally, we employ HumanEval \citep{codex} and MBPP \citep{mbpp}, both established benchmarks for code generation evaluation. Experimental evidence suggests that GRPO remains effective even when DeepSeekMath-Instruct 7B achieves high benchmark performance. Additionally, we offer an overarching characterization of the base model's competencies outside of mathematics, spanning natural language comprehension, logical reasoning, and software programming. Instead, GRPO leverages the mean reward of a set of sampled outputs for the same question as a baseline. As illustrated in Figure \ref{fig:rl-analysis}, the empirical results reveal that Online RFT achieves substantially better performance compared to RFT across two evaluated benchmarks. In particular, DeepSeekMath-Base sets a new standard on the competitive MATH dataset, exceeding the highest performing open-source base models by more than 10\% in absolute terms. We posit that \textbf{WEAK-TO-STRONG} alignment approaches \citep{burns2023weak} have the potential to fundamentally reshape the landscape of learning algorithms. This work presents DeepSeekMath~7B, an advancement built upon continued pre-training of DeepSeek-Coder-Base-v1.5 7B using a dataset of 120B math-focused tokens derived from Common Crawl, as well as accompanying natural language and programming code resources. \item \textbf{Multilingual}:
    The DeepSeekMath~Corpus is rich in linguistic diversity, with English and Chinese comprising the largest portions. Table \ref{tab:corpora_comparison} provides evidence that models trained with the DeepSeekMath~Corpus consistently achieve the highest scores. \section{Discussion}
This section elaborates on our insights gleaned from both pre-training and RL-based experiments. We undertake a thorough evaluation of DeepSeekMath-Base 7B’s mathematical proficiency, concentrating on three core tasks: generating fully self-contained mathematical answers unaided by external systems, leveraging computational tools for mathematical problem solving, and engaging in formal theorem proving. These results not only outperform all other open-source models within the 7B–70B category, but also exceed those of most closed-source alternatives. \end{itemize}

Given these open questions, additional, more nuanced studies are needed to further clarify arXiv’s impact. If a domain has more than 10\% of its pages collected, we designate it as math-oriented (e.g., \url{mathoverflow.net}). We outline potential research directions concerning each of these components. To mitigate the risk of contaminating downstream benchmarks, we implement filtering procedures as in \cite{deepseek-coder}. Additionally, rather than employing the form of KL penalty from (\ref{eq

\begin{algorithm}[t]
  \small
  \caption{Iterative Group Relative Policy Optimization}
  \textbf{Input} initial policy model $\pi_{\theta_{\text{init}}}$; reward models $r_\varphi$; task prompts $\mathcal{D}$; 
  hyperparameters $\varepsilon$, $\beta$, $\mu$
  \begin{algorithmic}[1]
    \State Initialize policy model $\pi_\theta \leftarrow \pi_{\theta_{\text{init}}}$
    \For{iteration = 1, \dots, I}
      \State Assign current policy to reference model: $\pi_{ref} \leftarrow \pi_{\theta}$
      \For{step = 1, \dots, M}
        \State Draw mini-batch $\mathcal{D}_b$ from $\mathcal{D}$
        \State Set $\pi_{\theta_{old}} \leftarrow \pi_{\theta}$ 
        \State For each question $q$ in $\mathcal{D}_b$, generate $G$ samples $\{o_i\}_{i=1}^G$ from $\pi_{\theta_{old}} (\cdot \mid q)$
        \State Compute associated rewards $\{r_i\}_{i=1}^{G}$ for outputs $o_i$ using $r_{\varphi}$
        \State Estimate group relative advantages $\hat{A}_{i,t}$ at each $t$-th token in $o_i$
        \For{GRPO iteration = 1, \dots, $\mu$}
          \State Update policy model $\pi_{\theta}$ to maximize the GRPO objective (Equation \ref{eq:GC-GRPO})
        \EndFor
      \EndFor 
      \State Periodically train reward model $r_\varphi$ with replay-based updates. The term $A_t$ indicates the advantage, which is estimated via Generalized Advantage Estimation (GAE) \citep{gae}, utilizing both future reward signals $\{r_{\ge t}\}$ and a parameterized value estimator $V_{\psi}$. These results provide a partial resolution to the frequently debated issue: \textit{does code training contribute to reasoning skills?} In the context of mathematical reasoning, our evidence suggests a positive impact. Furthermore, it surpasses Minerva 540B \citep{minerva}, a proprietary base model that is 77 times larger, is built upon PaLM \citep{palm}, and is subject to additional mathematical training. Leveraging a carefully engineered data curation workflow, we assemble the DeepSeekMath Corpus—a refined dataset comprising 120B tokens sourced from web pages specifically screened for mathematical relevance. Unless explicitly noted otherwise, all experimental setups follow the procedures described in Section \ref{sec:quality-policy}. Online sampling entails generating training data directly from the on-policy exploration undertaken by the model in real-time, whereas offline sampling utilizes data produced by the initial SFT model, independent of the evolving training policy. Broadly, for each method, the training gradient with respect to parameters $\theta$ may be expressed as:

\begin{equation}
    \nabla_{\theta}\mathcal{J}_{\textcolor{red}{\mathcal{A}}}(\theta) = \mathbb{E}[\underbrace{(q,o) \sim \textcolor{red}{\mathcal{D}}}_{Data \ Source}]\left( \frac{1}{|o|} \sum_{t=1}^{|o|}  \underbrace{GC_{{\mathcal{A}}}(q, o, t, \textcolor{red}{\pi_{{rf}}})}_{Gradient \ Coefficient}  \nabla_{\theta}\log \pi_{\theta}(o_t | q, o_{<t})\right). \end{itemize}

\subsection{Training and Evaluating DeepSeekMath-Base 7B}

Here, we present DeepSeekMath-Base 7B, a foundational model engineered for robust mathematical reasoning. \end{itemize}

\textbf{One-Stage Training}

\begin{itemize}[topsep=0pt]
    \item \textbf{Math Training for 150B Tokens}: The DeepSeek-LLM 1.3B model is exposed solely to 150B mathematics tokens. Additionally, we introduce a unified framework for interpreting diverse training techniques, categorizing them as variants or simplifications of RL. \begin{itemize}[topsep=0pt]
    \item \textbf{English mathematical datasets}: Our process includes annotating GSM8K and MATH datasets with solutions that utilize tools, along with incorporating a selected portion of MathInstruct \citep{MathInstruct} and the training split of Lila-OOD \citep{lila}, where questions are resolved via CoT or PoT methods. \item \textbf{Large-scale}:
    The scale of DeepSeekMath~Corpus significantly exceeds that of all comparable math corpora. In the early training epochs, Online RFT exhibits performance on par with RFT; however, it subsequently secures a notable performance edge, highlighting the benefits of online data collection. We aim to partially validate this claim within the scope of mathematical reasoning: namely, that code-focused training leads to improvements in mathematical problem-solving, independent of whether external tools are employed. \section{Reinforcement Learning}

\subsection{Group Relative Policy Optimization} Reinforcement learning (RL) methods have shown substantial promise in advancing the mathematical reasoning abilities of LLMs subsequent to the Supervised Fine-Tuning (SFT) phase \citep{wang2023math,wizardmath}. Furthermore, by evaluating 64 samples, DeepSeekMath~7B achieves a self-consistency score of 60.9\% on the MATH benchmark. The maximum learning rate is set to 5.3e-4. Drawing from \cite{wang2023math}, we further investigate process supervision, which evaluates and supplies rewards following every individual reasoning step. We identify three crucial research avenues for improving reward models: 1) \textbf{Strategies to improve the generalization capabilities of the reward model.} For reinforcement learning to lead to true advancements in large language models—rather than merely stabilizing their distribution—the reward model must be capable of generalizing to both out-of-distribution queries and sophisticated decoding scenarios; 2) \textbf{Mechanisms to account for reward model uncertainty.} Capturing uncertainty may serve as a conceptual bridge linking imperfect reward models and weak-to-strong learning strategies; 3) \textbf{Methods for efficiently constructing high-quality process reward models} that generate detailed, step-level feedback for reasoning tasks \citep{lightman2023let,wang2023math}. We also outline a cohesive framework to contextualize various techniques and identify promising avenues for advancing reinforcement learning. Our process removes any portion of text containing a 10-gram exact match to substrings from these benchmarks. \paragraph{Algorithms} Algorithms utilize both the data and the received reward signals to compute gradient coefficients, which are then employed to adjust model parameters. We also examine the impact of iterative RL, implementing two consecutive rounds in our study. For RL training, we draw on approximately 144K chain-of-thought questions drawn from GSM8K and MATH within the SFT dataset. Model performance is assessed on both producing complete natural language solutions without external tools, as well as solving the problems via Python. \item \textbf{Training on a mixture of 400B Code Tokens and 150B Math Tokens}: Since code pretraining followed by math fine-tuning adversely affects code generation capability, we investigate a regime where code and math tokens are presented together in a single-phase mixture to evaluate whether this strategy not only supports mathematical reasoning but also addresses catastrophic forgetting. \paragraph{Formal Mathematics}

The automation of formal proofs plays a critical role in ensuring both the accuracy and efficiency of mathematical verification, and has garnered significant attention recently. \end{itemize}

A summary of these methodologies is provided in Table \ref{tab:data_policy}. \item We report insights gathered from our mathematical training procedures. Training is performed using batches of 4M tokens, each with a context window of 4K tokens. If a benchmark sequence is shorter than 10 grams but at least 3 grams in length, we similarly exclude exact matches to ensure no overlap with evaluation data. \section{Conclusion, Limitation, and Future Work} In this work, we introduce DeepSeekMath, a model that surpasses all open-source alternatives on the MATH benchmark at the competition level, and comes close to matching the results of proprietary models. Our empirical analysis encompasses a range of experimental comparisons (e.g., online versus offline training, outcome versus process-based supervision, single-turn versus iterative RL), providing a comprehensive examination of key paradigm components. Results in Table \ref{tab:base_math_pot_proof} indicate that DeepSeekMath-Base 7B exceeds the performance of the previous state-of-the-art Llemma 34B. Surprisingly, our results indicate that exposure to arXiv documents does not substantively bolster mathematical reasoning proficiency. The composition of the training dataset is distributed as follows: 56\% originates from the DeepSeekMath Corpus, 4\% from AlgebraicStack, 10\% from arXiv, 20\% from code repositories on Github, with the last 10\% comprising natural language text from the Common Crawl dataset, incorporating both English and Chinese sources. \end{itemize}

\subsection{Summary of Evaluations and Metrics}

\begin{itemize}[topsep=0pt]
    \item \textbf{English and Chinese Mathematical Reasoning}: We perform thorough evaluations of our models using both English and Chinese datasets that encompass mathematics problems from elementary through university curricula. Additionally, as illustrated in Figure \ref{fig:corpora_comparisons}, DeepSeek-LLM 1.3B fine-tuned on DeepSeekMath outperforms its counterpart trained on Proof-Pile-2 at 50B tokens (which constitutes a single epoch for Proof-Pile-2), demonstrating that the average quality of DeepSeekMath surpasses that of competing corpora. This evaluation is conducted on miniF2F \citep{minif2f}, which serves as a benchmark for formal mathematics at the Olympiad level. Model optimization is conducted over 500 steps, employing a batch size of 256 and a fixed learning rate set at 5e-5. Finally, we elucidate how reinforcement learning contributes to the enhancement of instruction-tuned models and propose future research trajectories for realizing more effective RL techniques within this unified conceptual framework. The dataset size is further refined through preliminary pre-training experiments, considering the top 40B, 80B, 120B, and 160B tokens. The initial reward model is finetuned from DeepSeekMath-Base 7B, adopting a learning rate of 2e-5. Overall, DeepSeekMath-Base 7B significantly outperforms the general-purpose Mistral 7B \citep{mistral} on all three reasoning and coding tasks. \subsection{Lessons Learnt in Pre-Training}

We begin by detailing our observations related to the pre-training process. Following the initial round of data gathering, a substantial number of mathematical web pages remain uncaptured, predominantly due to the limited diversity within the positive training examples used for the fastText model. We categorize various prominent strategies in alignment with this unified schema:

\begin{itemize}
    \item  \textbf{Supervised Fine-tuning (SFT)}: SFT adapts a pretrained network using curated data annotated or chosen by human evaluators. Table \ref{tab:corpora_comparison} reveals that training with this corpus leads to marked improvements in mathematical reasoning for both English and Chinese tasks. Questions from other SFT categories are excluded to isolate the effects of RL on benchmark datasets absent from the training distribution. DeepSeekMath-Base is built upon DeepSeek-Coder-Base-v1.5 7B \citep{deepseek-coder}, as our findings indicate that initializing from a code-oriented pretrained model yields superior results relative to leveraging a generic LLM. \label{eq:objective}
\end{equation}

This framework highlights three fundamental elements: 1) \textit{Data Source $\mathcal{D}$}, specifying the set of samples for training; 2) \textit{Reward Function $\pi_{{rf}}$}, which generates the feedback signal utilized during optimization; 3) \textit{Algorithm $\mathcal{A}$}, responsible for processing input data and reward feedback to compute the gradient coefficient $GC$, thus modulating the strength of learning updates. In other words, \textbf{the observed gains arise from raising the correct response’s likelihood within TopK selections rather than from strengthening the model’s intrinsic abilities.} This observation is reminiscent of findings by \citep{wang2023making}, who reported a \textbf{misalignment problem} in reasoning tasks for SFT models, and further showed that alignment strategies focused on preference correction can enhance SFT reasoning performance \citep{yuan2023rrhf,song2023preference,wang2023making}. Each output is evaluated by a reward model, producing a reward set $\mathbf{r}=\{r_1, r_2, \cdots, r_G\}$. Adhering to the established practices for DeepSeek LLMs, we adopt the AdamW optimizer \citep{adamW} with $\beta_1=0.9$, $\beta_2=0.95$, and $\mathrm{weight\_decay}=0.1$. Moreover, GRPO+PS achieves better results than GRPO+OS, underscoring the advantage of employing step-aware, fine-grained gradient adjustments. Table \ref{tab:understanding_reasoning_code} demonstrates that DeepSeekMath-Base 7B delivers substantial improvements over its predecessor, DeepSeek-Coder-Base-v1.5 \citep{deepseek-coder}, on both MMLU and BBH, highlighting the beneficial effects of targeted mathematical training on broader linguistic comprehension and reasoning tasks. Moreover, due to its model size, DeepSeekMath~is outperformed by GPT-4 on few-shot learning benchmarks: while GPT-4 benefits noticeably from few-shot prompts, DeepSeekMath~displays comparable results in both zero-shot and few-shot evaluations. According to Table \ref{tab:base_math_pot_proof}, DeepSeekMath-Base 7B achieves strong results in automatic proof formalization. \subsection{Contributions}
We contribute by presenting large-scale math pre-training and investigating reinforcement learning techniques. On the challenging MATH dataset, our model outperforms every open-source competitor and most proprietary models (such as Inflection-2 and Gemini Pro) by a margin of at least 9\% absolute. Additionally, we observe that training with mathematical datasets not only boosts the model's mathematical skills but also enhances its performance on broader benchmarks, including MMLU \citep{mmlu} and BBH \citep{bbh}, thereby strengthening general reasoning abilities. For model training, 500,000 samples from the seed corpus serve as positive examples, while another 500,000 randomly chosen web pages from Common Crawl serve as negative samples. \item Leveraging this comprehensive paradigm, we delve into the mechanisms that make reinforcement learning effective and outline prospective directions for advancing the reinforcement learning of LLMs. This dataset is approximately 7-fold greater than the math corpus utilized by Minerva \citep{minerva}, and surpasses the newly introduced OpenWebMath \citep{openwebmath} by a factor of 9. For an expanded explanation and a detailed derivation, see Appendix \ref{app:analysis-rl}. Consequently, PPO entails joint training of both the policy and an auxiliary value function; in order to counteract excessive reward model exploitation, it is standard practice to impose a per-token KL penalty, as measured against a reference model, within the reward computation for each generated token \citep{ouyang2022training}, formulated as:

\begin{equation}
    r_{t} = r_\varphi(q, o_{\le t}) - \beta \log\frac{\pi_{\theta}(o_{t}|q, o_{<t})}{\pi_{ref}(o_{t}|q, o_{<t})},
\label{eq:PPO-reward}
\end{equation}

where $r_\varphi$ indicates the reward function, $\pi_{ref}$ designates the reference model (typically initialized from the SFT checkpoint), and $\beta$ specifies the magnitude of the KL penalty. The results indicate that DeepSeekMath-Base achieves notable effectiveness in few-shot autoformalization tasks. By contrast, other baseline corpora are considerably smaller, necessitating multiple passes through the dataset during training, which leads to model performance saturating rapidly. According to Equation \ref{eq:objective}, prevailing methods largely place complete \textbf{TRUST} in the feedback from the reward function to modulate the conditional probability associated with specific tokens, either promoting or suppressing them. Conducting code pre-training before embarking on mathematical training enhances the model's problem-solving competencies—regardless of tool usage. Specifically, for a given prompt $q$, a set of outputs $\{o_1, o_2, \cdots, o_G\}$ is generated from $\pi_{\theta_{old}}$, and the policy is improved by maximizing the following objective:

\begin{equation}
\footnotesize
\begin{split}
    \mathcal{J}_{GRPO}(\theta) &= \mathbb{E}{[q \sim P(Q), \{o_i\}_{i=1}^G \sim \pi_{\theta_{old}}(O|q)]}  \\
    & \frac{1}{G}\sum_{i=1}^G\frac{1}{|o_i|} \sum_{t=1}^{|o_i|} \left\{ \min \left[ \frac{\pi_\theta(o_{i,t} | q, o_{i,<t})}{\pi_{\theta_{old}}(o_{i,t} | q, o_{i,<t})} \hat{A}_{i,t}, \text{clip} \left( \frac{\pi_\theta(o_{i,t} | q, o_{i,<t})}{\pi_{\theta_{old}}(o_{i,t} | q, o_{i,<t})}, 1 - \varepsilon, 1 + \varepsilon \right)  \hat{A}_{i,t} \right] - \beta \mathbb{D}_{KL}\left[\pi_{\theta} || \pi_{ref}\right]\right\} ,
\end{split}
\label{eq:GRPO-obj}
\end{equation}

Here, the symbols $\varepsilon$ and $\beta$ remain as tunable hyper-parameters, while $\hat{A}_{i,t}$ is an advantage term evaluated using only the intra-group relative reward structure, as will be elaborated in subsequent subsections. Moreover, as the DeepSeekMath Corpus supports multiple languages, our experiments reveal gains on Chinese mathematical benchmarks \citep{wei2023cmath,agieval}. This substantial enhancement in mathematical reasoning is rooted in two main elements: firstly, the exploitation of vast quantities of openly available web material, systematically curated through an advanced data filtering pipeline; secondly, the development of Group Relative Policy Optimization (GRPO), an extension of Proximal Policy Optimization (PPO), which both elevates mathematical reasoning capabilities and streamlines PPO memory utilization. We posit that our methodology for assembling mathematical datasets sets a promising foundation for continued work in this area, acknowledging the considerable potential for future advancements. As illustrated in Figure \ref{fig:rl-analysis}, GRPO outperforms online RFT, thereby emphasizing the effectiveness of modifying positive and negative gradient coefficients. \item \textbf{Natural Language Understanding, Reasoning, and Code}: For a holistic evaluation of general linguistic understanding, reasoning ability, and coding proficiency, we test DeepSeekMath-Base on the Massive Multitask Language Understanding (MMLU) benchmark \citep{mmlu}, which includes 57 distinct multiple-choice subject areas, as well as on BIG-Bench Hard (BBH) \citep{bbh}, composed of 23 tasks primarily targeting complex multi-step reasoning. This paper presents DeepSeekMath, a specialized language model that demonstrates a marked improvement over open-source models in mathematical reasoning and approaches the capabilities of GPT-4 on standard academic benchmarks. Results depicted in Figure \ref{fig:iter-rl} reveal that iterative RL markedly boosts performance, particularly during the initial iteration. DeepSeekMath~builds on DeepSeek-Coder-v1.5 7B, receiving an additional 500B tokens of continued training, of which 120B math-related tokens are collected from Common Crawl. Furthermore, we offer a unified perspective for interpreting various approaches—including Rejection Sampling Fine-Tuning (RFT) \citep{yuan2023scaling}, Direct Preference Optimization (DPO) \citep{dpo}, PPO, and GRPO—by situating them within a shared RL paradigm characterized by either direct or simplified reinforcement learning methodologies. The maximum sequence length and training batch size are both set to 1024. \subsection{Training and Evaluating DeepSeekMath-RL}

Our reinforcement learning experiments utilize the DeepSeekMath-Instruct 7B model. To address this, we investigate the approach of iterative RL with GRPO. For each prompt, $64$ candidate responses are sampled. \subsection{Assessing the DeepSeekMath~Corpus Quality}

To evaluate the utility of the DeepSeekMath~Corpus, we conduct pre-training studies in comparison to several contemporary math corpora:

\begin{itemize}[topsep=0pt]
    \item \textbf{MathPile} \citep{mathpile}: a heterogeneous corpus (8.9B tokens) compiled from sources such as textbooks, Wikipedia, ProofWiki, CommonCrawl, StackExchange, and arXiv, with arXiv contributions making up more than 85\% of the data;
    \item \textbf{OpenWebMath} \citep{openwebmath}: mathematics-focused content curated from CommonCrawl, yielding a collection of 13.6B tokens;
    \item \textbf{Proof-Pile-2} \citep{llemma}: a mathematics corpus comprising OpenWebMath, AlgebraicStack (containing 10.3B tokens of math code), and arXiv papers (contributing 28.0B tokens). To ensure quality, pages are scored with the fastText model and only those receiving the highest scores are retained. According to Equation \ref{eq:objective}, this framework is structured around three main elements: Data Source, Algorithm, and Reward Function. \subsection{Insights of Reinforcement Learning}
\subsubsection{A Move Toward a Unified Framework} Here, we seek to establish a general framework for examining various training methodologies—including SFT, RFT, DPO, PPO, GRPO—and supplement this with experimental investigation into factors underpinning this paradigm. Model performance is assessed for mathematical reasoning, both with and without the use of external tools, across four quantitative reasoning benchmarks in both English and Chinese. Moreover, in the PPO paradigm, the value function primarily serves to provide a baseline to minimize variance in the computed advantage. The English datasets include GSM8K \citep{gsm8k}, MATH \citep{MATH}, SAT \citep{llemma}, OCW Courses \citep{minerva}, and MMLU-STEM \citep{mmlu}; for Chinese, we use MGSM-zh \citep{mgsm}, CMATH \citep{wei2023cmath}, Gaokao-MathCloze \citep{agieval}, and Gaokao-MathQA \citep{agieval}. Next, we carry out a manual labeling process, pinpointing URLs that specifically pertain to mathematical materials within these domains (e.g., \url{mathoverflow.net/questions}). Process supervision computes the advantage for every token as the cumulative sum of the future normalized rewards over subsequent reasoning steps, i.e., $\hat{A}_{i, t} = \sum_{index(j) \ge t} \widetilde{r}_i^{index(j)}$,
and optimizes the policy by maximizing the criterion detailed in equation

\subsubsection{Iterative RL with GRPO} During the reinforcement learning process, the reward model initially used may eventually fall short in effectively supervising the evolving policy model. The key observations are as follows: 1) DeepSeekMath-RL 7B achieves accuracy rates of 88.2\% on GSM8K and 51.7\% on MATH benchmarks when applying chain-of-thought methodologies. For this work, we only draw questions from the instruction tuning phase and rely on basic nucleus sampling for output generation. Here, we present Group Relative Policy Optimization (GRPO), our proposed RL approach designed to be both efficient and robust. The resulting rewards are standardized: each reward is centered by subtracting the batch mean, then scaled by dividing by the batch standard deviation. Furthermore, \textbf{efficient inference techniques} \citep{xia-etal-2023-speculative,leviathan2023fast,kwon2023efficient,xia2024unlocking}, which are critical for efficient policy model exploration, warrant significant attention. As detailed in Algorithm \ref{alg:iter-grpo}, this procedure involves regularly regenerating training datasets for the reward model by sampling from the latest policy model outputs. Additionally, code pretraining boosts mathematical reasoning abilities in settings without external tool use. \end{itemize}

\textbf{Exploration and Analysis of Reinforcement Learning}

\begin{itemize}[topsep=0pt]
    \item We present Group Relative Policy Optimization (GRPO), a novel reinforcement learning technique designed for both efficiency and effectiveness. 
\begin{abstract}
The intricacy and formal structure inherent in mathematical reasoning present formidable obstacles for language models. Training samples are concatenated at random until a maximum context size of 4K tokens is achieved. As detailed in Table \ref{tab:sft_rl_math}, when tool usage is not permitted during evaluation, DeepSeekMath-Instruct 7B achieves superior performance in stepwise mathematical reasoning. \subsubsection{Why RL Works?} In this study, we employ reinforcement learning using a selected portion of instruction tuning data, which substantially improves the performance over the instruction-tuned model. As shown in Figure \ref{fig:corpora_comparisons}, DeepSeek-LLM 1.3B, when trained on DeepSeekMath, achieves a notably sharper learning trajectory and more durable gains. For the initial iteration, only the top 40B tokens are preserved. The corpus is curated from Common Crawl (CC) utilizing a fastText-based classifier \citep{joulin2016fasttext}. The variables $q$ and $o$ represent the sampled questions and corresponding outputs, drawn from the question corpus and according to the distribution induced by the previous policy $\pi_{\theta_{old}}$. Although this represents a narrow scope in training data, DeepSeekMath-RL 7B nevertheless consistently surpasses DeepSeekMath-Instruct 7B in every evaluation metric, thereby highlighting the strengths of reinforcement learning. Instruction tuning and reinforcement learning further boost the resulting DeepSeekMath-Instruct and DeepSeekMath-RL models, which achieve over 50\% accuracy on the advanced MATH benchmark—a milestone not previously reached by any open-source model. The improved dataset then serves to refine the classifier, enhancing its subsequent selection accuracy. \item We establish a comprehensive framework to analyze various methodologies, such as RFT, DPO, PPO, and GRPO. We posit this limitation may account for our RL pipeline improving Maj@K but not Pass@K. Except for adjusting the peak learning rate to 4.2e-4 and utilizing a batch size of 10M tokens, we retain the training protocol detailed in Section \ref{sec:quality-policy}. Any mathematical pages associated with these URLs that were not retrieved previously are added to our set of positive examples, thereby enlarging the seed corpus. Here, GSM8K and MATH with chain-of-thought responses are categorized as in-domain, while all remaining benchmarks are treated as out-of-domain. In RL settings, it specifically denotes the unlabeled questions along with model-generated outputs. Utilizing only a subset of English-language instruction tuning samples, GRPO delivers marked improvements over the robust DeepSeekMath-Instruct, advancing both in-domain tasks (GSM8K: 82.9\% $\rightarrow$ 88.2\%, MATH: 46.8\% $\rightarrow$ 51.7\%) and out-of-domain tasks such as CMATH (84.6\% $\rightarrow$ 88.8\%) during the RL stage. DeepSeekMath-Base also excels on Chinese datasets, arguably due to the inclusion of high-quality multilingual data—contrary to earlier efforts \citep{minerva,llemma} that prioritized English-only pre-training. Each model is prompted to construct a Python script to solve given problems, with access to computational libraries like \textit{math} and \textit{sympy} for advanced calculations. Additionally, we observe gains in generalization to out-of-domain tasks throughout the reinforcement learning phase. To address these issues, as depicted in Figure \ref{fig:grpo}, we introduce Group Relative Policy Optimization (GRPO), which eliminates the need for training a separate value function, as is required in PPO. This strategy supplies the fastText model with a richer set of positive samples, allowing us to refine the model so it can recall a greater share of mathematics content in subsequent passes. Specifically, we exclude web pages that contain content from notable English math benchmarks (e.g., GSM8K \citep{gsm8k}, MATH \citep{MATH}) and Chinese benchmarks (such as CMATH \citep{wei2023cmath} and AGIEval \citep{agieval}). \subsubsection{Process Supervision RL with GRPO} Outcome-level feedback assigns a reward only at the end of an output, which may inadequately guide the learning process for complex reasoning problems. Nevertheless, the precise effects of arXiv pre-training on models' mathematical reasoning abilities remain underexplored. The optimization then proceeds by maximizing the objective specified in equation (\ref{eq:GRPO-obj}). \subsubsection{Code Training Benefits Mathematical Reasoning}

A widely held, though largely unsubstantiated, belief in the field is that training on code enhances reasoning ability. In subsequent stages, the divergence grows, leading to larger discrepancies in sampled data, which in turn amplifies the advantages of continual, real-time data acquisition. According to Table \ref{tab:code-math}, code pretraining alone, within the two-stage approach, already offers notable gains in tackling GSM8K and MATH datasets using Python tools, and further improvements accrue from subsequent math-specific tuning. \section{Math Pre-Training}

\subsection{Data Collection and Decontamination} 
This section details the methodology employed in assembling the DeepSeekMath Corpus using Common Crawl as the data source. For instance, even in PRM800K datasets \citep{lightman2023let}—which have been meticulously annotated by expert annotators—there persists an error rate of roughly 20\% in annotations\footnote{\url{https://github.com/openai/prm800k/issues/12\#issuecomment-1728491852}}. Beyond benchmark performance, these systems have shown remarkable utility in providing assistance for complex mathematical problem-solving tasks \citep{tao}. Nevertheless, leading models like GPT-4 \citep{gpt4} and Gemini-Ultra \citep{gemini} remain inaccessible to the public, while current open-source offerings substantially underperform relative to their proprietary counterparts. Through extensive experimentation—including studies of online versus offline optimization, outcome-based versus process-based supervision, and single-step versus iterative reinforcement learning—we systematically examine key factors underpinning this unified framework. \item While arXiv literature is a frequent component in math-oriented datasets, our analysis finds that its inclusion does not yield significant gains across the mathematical benchmarks considered here. This setup is used to assess whether code tokens offer unique benefits for mathematical reasoning compared to general-domain data. When models are allowed to leverage both natural language processing and programmatic tools for reasoning during evaluation, DeepSeekMath-Instruct 7B reaches nearly 60\% accuracy on MATH, surpassing all current open-source alternatives. Future efforts will focus on deploying our RL approach with out-of-distribution prompts and adopting \textbf{advanced sampling (decoding) strategies}, such as those leveraging tree-search \citep{yao2023tree}. Conversely, Online RFT does not exhibit such adaptive behavior; it reinforces all correct responses equivalently, omitting any penalization for incorrect outputs. 2) Significantly, the training regimen for DeepSeekMath-RL 7B is confined exclusively to instruction-tuning data in the chain-of-thought style for GSM8K and MATH, originating from DeepSeekMath-Instruct 7B. \end{document} \item Our DeepSeekMath-Base 7B pre-trained model attains results on par with the much larger Minerva 540B \citep{minerva}, suggesting that scaling model parameters is not the sole determinant of mathematical reasoning prowess. These datasets address a broad array of quantitative reasoning tasks (including GSM8K \citep{gsm8k}, MATH \citep{MATH}, and CMATH \citep{wei2023cmath}) as well as multiple-choice assessments (such as MMLU-STEM \citep{mmlu} and Gaokao-MathQA \citep{agieval}), spanning various mathematical disciplines and ranging in difficulty from basic arithmetic to advanced undergraduate mathematics. Moreover, with the inclusion of code-specific tokens in its continued training regimen, DeepSeekMath-Base 7B retains the strong code generation capabilities of DeepSeek-Coder-Base-v1.5 on both programming benchmarks.